% Nguồn SGK Toán 9 Tập hai — Chương VI, Luyện tập chung sau Bài 19, trang 18--20:
% https://cdn3.olm.vn/upload/thuvienso/4491669493-page-18-1771844540723.png
% https://cdn3.olm.vn/upload/thuvienso/4491669493-page-19-1771844541032.png
% https://cdn3.olm.vn/upload/thuvienso/4491669493-page-20-1771844541352.png

\section*{Luyện tập chung}

\begin{exercises}
    \textbf{Ví dụ 1} Cho hai hàm số: $y=\dfrac{3}{2}x^2$ và $y=-x^2$.

    \begin{enumerate}[label=\alph*)]
        \item Vẽ đồ thị của hai hàm số này trên cùng một mặt phẳng toạ độ.
        \item Tìm điểm $A$ thuộc đồ thị $y=\dfrac{3}{2}x^2$, điểm $B$ thuộc đồ thị $y=-x^2$, biết rằng $A$ và $B$ đều có hoành độ $x=-\dfrac{3}{2}$.
    \end{enumerate}

    \textbf{Giải}

    a) Lập bảng một số giá trị tương ứng giữa $x$ và $y$ của hai hàm số đã cho:

    \begin{center}
        \begin{tblr}{
            colspec={Q[c] *{5}{Q[c]}},
            hlines,
            vlines
        }
            $x$ & $-2$ & $-1$ & $0$ & $1$ & $2$ \\
            $y=\dfrac{3}{2}x^2$ & $6$ & $\dfrac{3}{2}$ & $0$ & $\dfrac{3}{2}$ & $6$
        \end{tblr}
        \qquad
        \begin{tblr}{
            colspec={Q[c] *{5}{Q[c]}},
            hlines,
            vlines
        }
            $x$ & $-2$ & $-1$ & $0$ & $1$ & $2$ \\
            $y=-x^2$ & $-4$ & $-1$ & $0$ & $-1$ & $-4$
        \end{tblr}
    \end{center}

    Từ đó vẽ được đồ thị của hai hàm số này như Hình 6.10:

    \begin{center}
        \begin{tikzpicture}[x=0.8cm,y=0.65cm]
            \coordinate (O) at (0,0);

            \draw[->] (-4.5,0) -- (4.5,0) node[below] {$x$};
            \draw[->] (0,-4.6) -- (0,6.8) node[left] {$y$};

            \foreach \x in {-4,-3,-2,-1,1,2,3,4}{
                \draw (\x,0.08) -- (\x,-0.08) node[below=2pt] {$\x$};
            }
            \foreach \y in {-4,-3,-2,-1,1,2,3,4,5,6}{
                \draw (0.08,\y) -- (-0.08,\y) node[left=2pt] {$\y$};
            }
            \node[below left=1pt of O] {$O$};

            \draw[dashed] (-2,6) -- (2,6);
            \draw[dashed] (-2,-4) -- (2,-4);
            \draw[dashed] (-2,-4) -- (-2,6);
            \draw[dashed] (2,-4) -- (2,6);

            \draw[red, domain=-2.12:2.12, smooth, samples=100] plot (\x,{1.5*\x*\x});
            \draw[blue, domain=-2.12:2.12, smooth, samples=100] plot (\x,{-(\x*\x)});

            \foreach \x/\y in {-2/6,-1/1.5,1/1.5,2/6}{
                \fill[red] (\x,\y) circle (1.2pt);
            }
            \foreach \x/\y in {-2/-4,-1/-1,1/-1,2/-4}{
                \fill[blue] (\x,\y) circle (1.2pt);
            }

            \node[rotate=70] at (0.95,4.1) {$y=\dfrac{3}{2}x^2$};
            \node[rotate=-70] at (1.05,-2.65) {$y=-x^2$};
        \end{tikzpicture}

        \textit{Hình 6.10}
    \end{center}

    b) Với $x=-\dfrac{3}{2}$, thay vào công thức $y=\dfrac{3}{2}x^2$ ta có: $y=\dfrac{3}{2}\cdot\left(-\dfrac{3}{2}\right)^2=\dfrac{27}{8}$. Vậy $A\left(-\dfrac{3}{2};\dfrac{27}{8}\right)$.

    Với $x=-\dfrac{3}{2}$, thay vào công thức $y=-x^2$ ta có: $y=-\left(-\dfrac{3}{2}\right)^2=-\dfrac{9}{4}$. Vậy $B\left(-\dfrac{3}{2};-\dfrac{9}{4}\right)$.

    \textbf{Ví dụ 2}

    Giả sử số lượng cá trong một hồ nào đó tăng, giảm theo công thức:
    \[
        P=50(100+15t-t^2),\ 0\le t\le 15.
    \]
    Ở đây $P$ là số lượng cá trong hồ sau $t$ năm tính từ ngày 01 tháng 01 năm 2020, khi lần đầu tiên số lượng cá trong hồ được ước tính. Hỏi theo mô hình này, thì:
    \begin{enumerate}[label=\alph*)]
        \item Vào thời điểm nào, số lượng cá trong hồ sẽ là $7\,500$ con?
        \item Vào thời điểm nào, số lượng cá trong hồ sẽ trở lại như tại thời điểm ban đầu vào ngày 01 tháng 01 năm 2020?
    \end{enumerate}

    \textbf{Giải}

    a) Ta cần tìm $t$ sao cho:
    \[
        P=7\,500,\text{ tức là }50(100+15t-t^2)=7\,500,\text{ hay }t^2-15t+50=0.
    \]
    Giải phương trình bậc hai này ta được hai nghiệm là $t=5$ hoặc $t=10$. Cả hai nghiệm này đều thích hợp.

    Giá trị $t=5$ ứng với thời điểm vào ngày 01-01-2025.

    Giá trị $t=10$ ứng với thời điểm vào ngày 01-01-2030.

    Vậy theo mô hình đã cho, có hai thời điểm mà số cá trong hồ là $7\,500$ con, đó là vào ngày 01-01-2025 hoặc vào ngày 01-01-2030.

    b) Số lượng cá trong hồ tại thời điểm ban đầu (ngày 01-01-2020) là:
    \[
        P(0)=50\cdot 100=5\,000\text{ (con)}.
    \]
    Ta cần tìm $t$ sao cho: $P=5\,000$, hay $15t-t^2=0$, tức là $t=0$ hoặc $t=15$.

    Giá trị $t=0$ ứng với thời điểm ban đầu, tức là ngày 01-01-2020.

    Giá trị $t=15$ ứng với thời điểm ngày 01-01-2035.

    Vậy theo mô hình đã cho, thì vào ngày 01-01-2035, số lượng cá sẽ trở lại như tại thời điểm ban đầu.

    \subsection{Bài tập}

    \begin{questions}[resume=SGK]
        \item Biết rằng parabol $y=ax^2$ ($a\ne 0$) đi qua điểm $A(2;4\sqrt{3})$.
        \begin{questions}
            \item Tìm hệ số $a$ và vẽ đồ thị của hàm số $y=ax^2$ với $a$ vừa tìm được.

            \answerline
            \answerline
            \answerline
            \answerline

            \item Tìm tung độ của điểm thuộc parabol có hoành độ $x=-1$.

            \answerline
            \answerline

            \item Tìm các điểm thuộc parabol có tung độ $y=5\sqrt{3}$.

            \answerline
            \answerline
        \end{questions}

        \item Công thức $E=\dfrac{1}{2}mv^2$ (J) được dùng để tính động năng của một vật có khối lượng $m$ (kg) khi chuyển động với vận tốc $v$ (m/s) (theo \textit{Vật lí đại cương, NXB Giáo dục Việt Nam, 2016}).
        \begin{questions}
            \item Giả sử một quả bóng có khối lượng 2 kg đang bay với vận tốc 6 m/s. Tính động năng của quả bóng đó.

            \answerline
            \answerline

            \item Giả sử động năng của quả bóng đang bay có khối lượng 1,5 kg là 48 J, hãy tính vận tốc bay của quả bóng đó.

            \answerline
            \answerline
        \end{questions}

        \item Cho hình chóp tam giác đều có đáy là tam giác đều cạnh $a$ (cm) và chiều cao 10 cm.
        \begin{questions}
            \item Tính diện tích đáy $S$ của hình chóp theo $a$.

            \answerline
            \answerline

            \item Từ kết quả ở câu a, tính thể tích $V$ của hình chóp theo $a$ và tính giá trị của $V$ khi $a=4$ cm.

            \answerline
            \answerline

            \item Nếu độ dài cạnh đáy giảm đi hai lần thì thể tích hình chóp thay đổi thế nào?

            \answerline
            \answerline
        \end{questions}

        \item Sử dụng công thức nghiệm hoặc công thức nghiệm thu gọn, giải các phương trình sau:
        \begin{questions}
            \item $x^2-2\sqrt{5}x+1=0$;

            \answerline
            \answerline

            \item $3x^2-9x+3=0$;

            \answerline
            \answerline

            \item $11x^2-13x+5=0$;

            \answerline
            \answerline

            \item $2x^2+2\sqrt{6}x+3=0$.

            \answerline
            \answerline
        \end{questions}

        \item Sử dụng máy tính cầm tay, tìm nghiệm gần đúng của các phương trình sau (làm tròn kết quả đến chữ số thập phân thứ hai):
        \begin{questions}
            \item $\sqrt{2}x^2-\sqrt{5}x-1=0$;

            \answerline
            \answerline

            \item $x^2-(\sqrt{3}-1)x-\sqrt{7}=0$.

            \answerline
            \answerline
        \end{questions}

        \item Từ một tấm tôn hình vuông, người ta cắt bỏ bốn hình vuông có độ dài cạnh 8 cm ở bốn góc, sau đó gập thành một chiếc thùng có dạng hình hộp chữ nhật không có nắp và có thể tích bằng 200 cm$^3$. Tính độ dài cạnh của tấm tôn hình vuông ban đầu.

        \answerline
        \answerline
        \answerline

        \item Giả sử doanh thu (nghìn đồng) của một cửa hàng bán phở trong một ngày có thể mô hình hoá bằng công thức $R(x)=x(220-4x)$ với $30\le x\le 50$, trong đó $x$ (nghìn đồng) là giá tiền của một bát phở. Nếu muốn doanh thu trong ngày của cửa hàng đạt 3 triệu đồng thì giá bán của mỗi bát phở phải là bao nhiêu?

        \answerline
        \answerline
        \answerline
    \end{questions}
\end{exercises}
